\input{sol_pre.tex}
\input{q_2024-04-28.tex}
\begin{document}
\input{sol_head.tex}

\Savol{1}
$2x+3y=6z$ $\Rightarrow$ $2x=3(2z-y)$. $\gcd(2;3)=1$ bölgani uçun $x$ soni $3$ ga bölinadi.
\Javob{B}

\Savol{2}
Içkaridan boşlaymiz: $1-\frac14=\frac34$, \ $\left(\frac34\right)^{-1}=\frac43$, \ $1-\frac43=-\frac13$, \ $\left(-\frac13\right)^{-3}=-27$,\\
$1-(-27)=28$, \ $28^{-1}=\frac1{28}$, \ $1-\frac1{28}=\frac{27}{28}$, \ $\left(\frac{27}{28}\right)^{-1}=\dfrac{28}{27}$.
\Javob{C}

\Savol{3}
Oltin $g$, kumuş $s$ gramm: \ $g+s=1060$, \ $\dfrac g{19}+\dfrac s{10}=70$.\\
Tekşiramiz: $g=760$, $s=300$: \ $\dfrac{760}{19}+\dfrac{300}{10}=40+30=70$.
\Javob{A}

\Savol{4}
$10\,000\cdot1{,}15^2=10\,000\cdot1{,}3225=13\,225$.
\Javob{D}

\Savol{5}
$3^{8-2}\cdot2^{-4+8}=2^4\cdot3^6$ $\Rightarrow$ $m=4$, $n=6$, \ $m+n=10$.
\Javob{C}

\Savol{6}
$\sqrt{4-2\sqrt3}=\sqrt3-1$; \ $22-30\left(\sqrt3-1\right)=52-30\sqrt3=\left(3\sqrt3-5\right)^2$ \ ($3\sqrt3>5$):
\[3\sqrt3-5+5+3\sqrt3=6\sqrt3.\]
\Javob{B}

\Savol{7}
$t=\sqrt[3]2$, \ $t^3=2$: \ $\dfrac3{2t(1+t)}-\dfrac3{2t(1-t)}=\dfrac3{2t}\cdot\dfrac{-2t}{1-t^2}=\dfrac3{t^2-1}$.
\[\frac3{t^2-1}=\frac{3\left(t^4+t^2+1\right)}{t^6-1}=t^4+t^2+1;\qquad t^4+t^2+1+t^2=\left(t^2+1\right)^2=\left(\sqrt[3]4+1\right)^2.\]
\Javob{C}

\Savol{8}
\[\frac{(a-2b)(a+2b)-5(a-2b)}{(a+2b-5)(a+2b+5)}=\frac{a-2b}{a+2b+5}=\frac{0{,}8}{8}=0{,}1.\]
\Javob{B}

\Savol{9}
Kamayuvçi: $a_5<a_3=1$ $\Rightarrow$ $a_5=-5$; \ $2d=-6$, $d=-3$, $a_1=a_3-2d=7$. \ $a_1d=-21$.
\Javob{B}

\Savol{10}
$b_1=2$, \ $q=\dfrac23$: \ $S=\dfrac2{1-\frac23}=6$.
\Javob{C}

\Savol{11}
$x\ge2$ va $x\le1$ bir vaqtda bajarilmaydi — aniqlaniş sohasi böş.
\Javob{D}

\Savol{12}
$\sqrt5-2=\left(\sqrt5+2\right)^{-1}$, asos $\sqrt5+2>1$:
\[x-1\le-\frac{x-1}{x+1}\ \Leftrightarrow\ \frac{(x-1)(x+2)}{x+1}\le0.\]
$x>1$ da ifoda musbat, $x=1$ da $0$. Musbat butun yeçim faqat $x=1$.
\Javob{A}

\Savol{13}
$t=\ln x$: \ $2^t=1+\left(\sqrt3\right)^t$ $\Leftrightarrow$ $\left(\frac12\right)^t+\left(\frac{\sqrt3}2\right)^t=1$.\\
Çap tomon kamayuvçi, $t=2$ da $\frac14+\frac34=1$. Yagona ildiz $x=e^2$.
\Javob{B}

\Savol{14}
Tengsizlik faqat $x^2-x-6=0$, ya'ni $x=3$, $x=-2$ da bajarilmaydi. Natural son — faqat $3$.
\Javob{B}

\Savol{15}
$405=3^4\cdot5$, \ $\frac{25}3=5^2\cdot3^{-1}$:
\[c=\frac{\log_2405}{\log_2\frac{25}3}=\frac{4a+b}{2b-a}.\]
\Javob{D}

\Savol{16}
$\sin\left(\frac\pi4-x\right)=\frac{\cos x-\sin x}{\sqrt2}$ $\Rightarrow$ $\sin^2\left(\frac\pi4-x\right)=\frac{(\sin x-\cos x)^2}2$. Ifoda $2$.
\Javob{A}

\Savol{17}
$2\sin\left(2x-\frac\pi6\right)=\sqrt3$:
\[2x-\frac\pi6=\frac\pi3+2\pi k\ \Rightarrow\ x=\frac\pi4+\pi k;\qquad 2x-\frac\pi6=\frac{2\pi}3+2\pi k\ \Rightarrow\ x=\frac{5\pi}{12}+\pi k.\]
$(0;5\pi)$ da har seriyadan $k=0,1,2,3,4$ — jami $10$ ta.
\Javob{C}

\Savol{18}
$t=|x|$: \ $-2\ge(t-2)(t+1)=t^2-t-2$ $\Rightarrow$ $t^2-t\le0$ $\Rightarrow$ $0\le t\le1$. Butun $x$: $-1,\ 0,\ 1$.
\Javob{D}

\Savol{19}
$x=\frac{1+\sqrt{13}}6$ — $3x^2-x-1=0$ ning ildizi. \ $g(x)=x\left(3x^2-x-1\right)+4=4$.
\Javob{C}

\Savol{20}
$g(x)>0$ ($x<0$ da) va $g(x)<0$ ($x>0$ da). $f'$ işorasi — $f$ ösişi yoki kamayişi:\\
$x_1$: $f\nearrow$, $g>0$ — ha; \ $x_2$: $f\searrow$, $g>0$ — yöq; \ $x_3$: $f\nearrow$, $g>0$ — ha;\\
$x_4$: $f\searrow$, $g<0$ — ha; \ $x_5$: $f\nearrow$, $g<0$ — yöq; \ $x_6$: $f\searrow$, $g<0$ — ha. \ Jami $4$ ta.
\Javob{C}

\Savol{21}
$\int x\sin x\,dx=-x\cos x+\sin x$:
\[\left(\frac{2\pi}3\cdot\frac12+\frac{\sqrt3}2\right)-\left(-\frac\pi3\cdot\frac12+\frac{\sqrt3}2\right)=\frac\pi3+\frac\pi6=\frac\pi2.\]
\Javob{B}

\Savol{22}
$t=\sqrt[4]{3x-2}\ge0$: \ $t^2+15t-16=0$ $\Rightarrow$ $t=1$ ($t=-16$ çet). \ $3x-2=1$, $x=1$. Köpaytma $1$.
\Javob{A}

\Savol{23}
$f(1)<0$, $f(2)>0$, $a>0$. Şoxlari yuqoriga, $f(1)<0$ $\Rightarrow$ $x_1<1<x_2$; $f(2)>0$ va $f(1)<0$ $\Rightarrow$ $x_2<2$.
\Javob{A}

\Savol{24}
$\dfrac3{12}\cdot360^\circ=90^\circ$.
\Javob{A}

\Savol{25}
Tomon $s$: $GH=s$, \ $HC=s+2\cdot\frac s{\sqrt2}=s\left(1+\sqrt2\right)$. \ $S_{GHCD}=s^2\left(1+\sqrt2\right)=12$.
\[S_8=2\left(1+\sqrt2\right)s^2=24.\]
\Javob{B}

\Savol{26}
Uçburçak tengsizligi: $x+5>3x-5$ $\Rightarrow$ $x<5$; \ $x+3x-5>5$ $\Rightarrow$ $x>2{,}5$. Eng katta natural $x=4$.
\Javob{B}

\Savol{27}
Kesuvçilar: $CF\cdot CD=CG\cdot CE=7\cdot12=84$ $\Rightarrow$ $CD=14$.\\
$\triangle CFG\sim\triangle CED$ ($\angle C$ umumiy, $\frac{CF}{CE}=\frac{CG}{CD}$):
\[\frac{FG}{ED}=\frac{CG}{CD}=\frac7{14}\ \Rightarrow\ FG=3.\]
\Javob{B}

\Savol{28}
\fig{\begin{tikzpicture}[scale=0.6,font=\small]
\coordinate (B) at (0,0); \coordinate (C) at (0,2); \coordinate (A) at (4,0); \coordinate (P) at (-2,0); \coordinate (Q) at (0,-4); \coordinate (M) at (2,1); \coordinate (N) at (-1.6,-0.8);
\draw (B)--(C)--(-2,2)--(P)--cycle; \draw (B)--(A)--(4,-4)--(Q)--cycle; \draw (A)--(C); \draw (P)--(Q); \draw[axs] (M)--(N);
\node[above] at (C) {$C$}; \node[above left] at (B) {$B$}; \node[right] at (A) {$A$}; \node[above right] at (M) {$M$}; \node[left] at (N) {$N$}; \node[above left] at (P) {$P$}; \node[below] at (Q) {$Q$};
\end{tikzpicture}}
$BC=2$, $BA=4$, $AC=2\sqrt5$, \ $BM=\frac{AC}2=\sqrt5$. \ $BP=BC$, $BQ=BA$, $\angle PBQ=90^\circ$ $\Rightarrow$ $\triangle PBQ=\triangle CBA$ ($B$ atrofida $90^\circ$ buriş). Şuning uçun $BM$ mediana davomi $PQ$ ga perpendikular: $BN$ — $\triangle PBQ$ balandligi.
\[BN=\frac{2\cdot4}{2\sqrt5}=\frac4{\sqrt5},\qquad MN=\sqrt5+\frac4{\sqrt5}=\frac9{\sqrt5}.\]
\Javob{C}

\Savol{29}
$\overrightarrow{FE}+\overrightarrow{ED}=\overrightarrow{FD}$, \ $\overrightarrow{AB}=\overrightarrow{ED}$. \ $|FD|=4\sqrt3$, $\overrightarrow{FD}$ va $\overrightarrow{ED}$ orasidagi burçak $30^\circ$:
\[\left|\overrightarrow{ED}+\overrightarrow{FD}\right|^2=16+48+2\cdot4\cdot4\sqrt3\cdot\frac{\sqrt3}2=112\ \Rightarrow\ 4\sqrt7.\]
\Javob{D}

\Savol{30}
Tekislikdan gipotenuzagaça masofa — $A$ va $B$ ning tekislikdan uzoqligi $12$. Katetlar: $\sqrt{16^2+12^2}=20$, \ $\sqrt{9^2+12^2}=15$. \ $AB=\sqrt{400+225}=25$.
\Javob{A}

\Savol{31}
$A\cap B=\{a,b,c,m\}$: \ $2^4-1=15$.
\Javob{C}

\Savol{32}
$4$ ta sondan takrorlanişli uçtalik: $\mathrm C_6^3=20$ ta. Eng yomon holat $8+8>14$ — hammasi uçburçak. Muntazam: $4$ ta. \ $P=\dfrac4{20}=\dfrac15$.
\Javob{A}

\newpage
\Savol{33}
\fig{\begin{tikzpicture}[scale=0.6,font=\small]
\draw (-1,0) rectangle (1,4); \fill[okgreen!30] (0,2) circle (1); \draw (0,2) circle (1);
\draw[axs] (0,2)--(1,2); \node[aux,above] at (0.5,2) {\scriptsize$r$};
\node[below] at (0,0) {$d=2r$}; \node[right] at (1,3.4) {$h=4r$};
\begin{scope}[shift={(5,2)}]
\draw (0,0) circle (1.3); \fill[okgreen!30] (45:1.3)--(135:1.3)--(225:1.3)--(315:1.3)--cycle; \draw (45:1.3)--(135:1.3)--(225:1.3)--(315:1.3)--cycle;
\draw[axs] (0,0)--(45:1.3); \node[aux,above left] at (45:0.6) {\scriptsize$r$}; \node[below] at (0,-1.3) {kesim};
\end{scope}
\end{tikzpicture}}
Göla radiusi $r$, balandligi $h=2d=4r$: \ $V=\pi r^2\cdot4r=12r^3$.

\textbf{33.} Eng katta şar radiusi $r$: $V_{\text{ş}}=\frac43\pi r^3=4r^3$. \ $\dfrac{4r^3}{12r^3}=\dfrac13$ — \textbf{C}.

\textbf{34.} Tösin kesimi — doiraga içki çizilgan kvadrat, tomoni $r\sqrt2$, yuzi $2r^2$: $V_{\text{t}}=2r^2\cdot4r=8r^3$. Çiqindi $12r^3-8r^3=4r^3$: \ $\dfrac{4r^3}{8r^3}=\dfrac12$ — \textbf{A}.

\textbf{35.} Konus hajmi gölaning $\frac13$ qismi, çiqindi $\dfrac23$ — \textbf{D}.
\Javob{33 — C, \ 34 — A, \ 35 — D}

\newpage
\Savol{36}
\fig{\begin{tikzpicture}[scale=0.32,font=\small]
\draw[-{Stealth[length=4pt]}] (-10.5,0)--(16,0) node[right] {$x$}; \draw[-{Stealth[length=4pt]}] (0,-8.5)--(0,9) node[left] {$y$};
\draw[thick] (6,4) circle (3); \draw[thick] (-6,4) circle (3);
\draw[okgreen,thick] (3,0) circle (7); \draw[axs] (3,0)--(6,4); \draw[axs] (3,0)--(-6,4);
\fill (3,0) circle (5pt); \fill (6,4) circle (5pt); \fill (-6,4) circle (5pt);
\node[below] at (3,0) {$(3;0)$}; \node[above] at (6,4) {$O_1$}; \node[above] at (-6,4) {$O_2$};
\node[aux] at (5.2,1.4) {$5$}; \node[aux] at (-1.5,2.8) {$\sqrt{97}$};
\end{tikzpicture}}
Birinçi tenglama: markazi $O_1(6;4)$ va $O_2(-6;4)$, radiusi $3$ bölgan ikki aylana ($x\ge0$ va $x\le0$ qismlarda töliq yotadi). Ikkinçisi — markazi $(3;0)$, radiusi $a$ bölgan aylana.\\
Markazlargaça masofa: $5$ va $\sqrt{97}\approx9{,}85$.\\
$O_1$ aylana bilan $2$ ta umumiy nuqta: $2<a<8$; \quad $O_2$ aylana bilan: $\sqrt{97}-3<a<\sqrt{97}+3$, ya'ni $6{,}85<a<12{,}85$.

a) $4$ ta yeçim: ikkala oraliq kesişmasi $6{,}85<a<8$ $\Rightarrow$ $a=7$.

b) $2$ ta yeçim: $2<a<\sqrt{97}-3$ yoki $8<a<\sqrt{97}+3$. Eng katta natural $a=12$.
\Javob{a) 7; \quad b) 12}

\newpage
\Savol{37}
$\frac12(\cos2x-\cos4x)=\sin3x\sin x$. Tengsizlik:
\[\sin3x\sin x+\sin3x-\cos3x-\cos3x\sin x\ge0\ \Leftrightarrow\ (1+\sin x)\left(\sin3x-\cos3x\right)\ge0.\]
$1+\sin x\ge0$ har doim; $1+\sin x=0$ da ($x=-\frac\pi2+2\pi k$) tengsizlik bajariladi. Qolgan hollarda
\[\sqrt2\sin\left(3x-\frac\pi4\right)\ge0\ \Rightarrow\ \frac\pi{12}+\frac{2\pi k}3\le x\le\frac{5\pi}{12}+\frac{2\pi k}3.\]

a) Eng kiçik musbat yeçim: $k=0$, $x=\dfrac\pi{12}$.

b) $k=-1$: $\left[-\frac{7\pi}{12};-\frac\pi4\right]$ — eng katta manfiy $-\frac\pi4$ ($-\frac\pi2$ undan kiçik).
\[\frac\pi{12}-\frac\pi4=-\frac\pi6.\]
\Javob{a) $\dfrac\pi{12}$; \quad b) $-\dfrac\pi6$}

\newpage
\Savol{38}
$f(x)+xf'(x)=\left(xf(x)\right)'$:
\[\int_a^b\left(xf(x)\right)'dx=bf(b)-af(a)=3b-2a=4.\]
$a+b=3$ bilan: $3(3-a)-2a=4$ $\Rightarrow$ $a=1$, $b=2$.

a) $ab=2$.

b) Çegaralar $a-4=-3$, $b+1=3$. \ $\dfrac{\sin3x}{x^2+4}$ — toq funksiya, simmetrik oraliqda integral $0$.
\Javob{a) 2; \quad b) 0}

\newpage
\Savol{39}
\fig{\begin{tikzpicture}[x=0.4cm,y=0.7cm,font=\small]
\draw[-{Stealth[length=4pt]}] (-9,0)--(15.3,0) node[right] {$x$}; \draw[-{Stealth[length=4pt]}] (0,-1.9)--(0,2.4) node[left] {$y$};
\draw[thick] plot[smooth,tension=0.55] coordinates {(-8,-1.0) (-7.3,0) (-5.8,1.6) (-4.6,0) (-3.9,-0.7) (-3.2,0) (-2.4,1.1) (-1.6,0) (-0.4,-1.4) (0.9,0) (2.2,1.8) (3.6,0) (4.9,-0.8) (6.2,0) (8.0,1.7) (9.6,0.9) (11.6,1.2) (14,0.25)};
\draw[fill=white] (-8,-1.0) circle (2.5pt); \draw[fill=white] (14,0.25) circle (2.5pt);
\foreach \x/\y in {-5.8/1.6,-3.9/-0.7,-2.4/1.1,-0.4/-1.4,2.2/1.8,4.9/-0.8,8.0/1.7,9.6/0.9,11.6/1.2} \fill[aux] (\x,\y) circle (2.5pt);
\foreach \x in {-7.3,-4.6,-3.2,-1.6,0.9,3.6,6.2} \fill[okgreen] (\x,0) circle (2.5pt);
\end{tikzpicture}}
a) $f'(x)=0$ — ekstremum nuqtalar (qizil): $5$ ta maksimum, $4$ ta minimum — jami $9$ ta.

b) $f(x)\cdot f'(x)=0$ $\Leftrightarrow$ $f(x)=0$ yoki $f'(x)=0$. \ $f(x)=0$: grafik $Ox$ ni $7$ marta kesadi (yaşil). Ekstremumlarning heç biri $Ox$ da emas:
\[9+7=16.\]
\Javob{a) 9 ta; \quad b) 16 ta}

\newpage
\Savol{40}
$x=1$: \ $2g(1)+3g(2)=1$; \qquad $x=2$: \ $2g(2)+3g(1)=2$.\\
Birinçisini $2$ ga, ikkinçisini $3$ ga köpaytirib ayiramiz:
\[9g(1)-4g(1)=6-2\ \Rightarrow\ g(1)=\frac45,\qquad g(2)=\frac{1-\frac85}3=-\frac15.\]
\Javob{a) $\dfrac45$; \quad b) $-\dfrac15$}

\newpage
\Savol{41}
\fig{\begin{tikzpicture}[scale=0.2,font=\small]
\coordinate (C) at (0,0); \coordinate (B) at (16,0); \coordinate (A) at (0,20.49); \coordinate (D) at (8,0); \coordinate (E) at (0,10.25);
\draw (A)--(B)--(C)--cycle; \draw[axs] (A)--(D); \draw[axs] (B)--(E); \draw (0,1.2)--(1.2,1.2)--(1.2,0);
\node[above] at (A) {$A$}; \node[right] at (B) {$B$}; \node[below left] at (C) {$C$}; \node[below] at (D) {$D$}; \node[left] at (E) {$E$};
\node[below] at (4,0) {$\frac a2$}; \node[below] at (12,0) {$\frac a2$}; \node[left] at (0,5) {$\frac b2$}; \node[left] at (0,15.4) {$\frac b2$};
\end{tikzpicture}}
$BC=a$, $AC=b$. \ $\triangle ACD$ va $\triangle BCE$ töğri burçakli:
\[b^2+\frac{a^2}4=22^2=484,\qquad a^2+\frac{b^2}4=19^2=361.\]

a) Qöşamiz: $\dfrac54\left(a^2+b^2\right)=845$ $\Rightarrow$ $AB^2=a^2+b^2=676$, \ $AB=26$.

b) Ayiramiz: $\dfrac34\left(b^2-a^2\right)=123$ $\Rightarrow$ $b^2-a^2=164$. \ $b^2=420$, $a^2=256$:
\[a=16,\quad b=2\sqrt{105},\qquad S=\frac12\cdot16\cdot2\sqrt{105}=16\sqrt{105}.\]
\Javob{a) 26; \quad b) $16\sqrt{105}$}

\newpage
\Savol{42}
\fig{\begin{tikzpicture}[scale=1.5,font=\small]
\foreach \n/\a in {A/112.5,B/67.5,C/22.5,D/-22.5,E/-67.5,F/-112.5,G/-157.5,H/157.5} {\coordinate (\n) at (\a:1); \node at (\a:1.15) {$\n$};}
\fill[okgreen!30] (A)--(D)--(E)--cycle;
\draw (A)--(B)--(C)--(D)--(E)--(F)--(G)--(H)--cycle; \draw (A)--(C); \draw (A)--(D); \draw (A)--(E);
\draw[ax] (0,0) circle (1); \fill (0,0) circle (1pt); \node[right] at (0,0.05) {$O$};
\draw ($(D)!0.12!(A)$)--($(D)!0.12!(A)!0.12!90:(A)$)--($(D)!0.12!(E)$);
\end{tikzpicture}}
Muntazam sakkizburçakning içki burçagi $135^\circ$.

a) $\triangle ABC$ da kosinuslar teoremasi:
\[AC^2=2a^2-2a^2\cos135^\circ=a^2\left(2+\sqrt2\right)=\left(4-2\sqrt2\right)\left(2+\sqrt2\right)=4\ \Rightarrow\ AC=2.\]

b) Taşqi aylana radiusi: $\angle AOC=90^\circ$ $\Rightarrow$ $AC=R\sqrt2$, $R=\sqrt2$.\\
$A$ va $E$ orasida $4$ ta tomon — $AE$ diametr: $AE=2\sqrt2$, $\angle ADE=90^\circ$ (diametrga tiralgan).
\[AD^2=AE^2-DE^2=8-\left(4-2\sqrt2\right)=4+2\sqrt2.\]
\[S_{ADE}=\frac12\,AD\cdot DE=\frac12\sqrt{\left(4+2\sqrt2\right)\left(4-2\sqrt2\right)}=\frac12\sqrt8=\sqrt2.\]
\Javob{a) 2; \quad b) $\sqrt2$}

\newpage
\Savol{43}
\fig{\begin{tikzpicture}[scale=0.8,font=\small]
\coordinate (A) at (0,0); \coordinate (E) at (2,0); \coordinate (B) at (4,0); \coordinate (D) at (2,3); \coordinate (C) at (6,3); \coordinate (F) at (2,1.5); \coordinate (G) at (2,1);
\fill[okgreen!30] (C)--(B)--(E)--(F)--cycle;
\draw (A)--(B)--(C)--(D)--cycle; \draw (D)--(E); \draw (A)--(C); \draw[axs] (F)--(C);
\draw (2,0.25)--(2.25,0.25)--(2.25,0);
\node[below left] at (A) {$A$}; \node[below] at (E) {$E$}; \node[below right] at (B) {$B$}; \node[above] at (D) {$D$}; \node[above right] at (C) {$C$}; \node[left] at (F) {$F$}; \node[left] at (G) {$G$};
\node[below] at (1,0) {$e$}; \node[below] at (3,0) {$e$};
\end{tikzpicture}}
$AE=EB=e$, $DE=h$: \ $S=2e\cdot h=48$, \ $eh=24$.

a) $S_{ADE}=\frac12eh=12$ $\Rightarrow$ $S_{EBCD}=48-12=36$.\\
$DF=\frac h2$, $DF\perp DC$, $DC=2e$: \ $S_{DFC}=\frac12\cdot\frac h2\cdot2e=\frac{eh}2=12$.
\[S_{CBEF}=36-12=24.\]

b) $AE\parallel DC$ $\Rightarrow$ $\triangle AGE\sim\triangle CGD$, \ $\frac{EG}{GD}=\frac{AE}{DC}=\frac12$ $\Rightarrow$ $EG=\frac h3$, \ $FG=\frac h2-\frac h3=\frac h6$.\\
$\triangle CFG$ da $FG$ asos, balandlik — $C$ dan $DE$ gaça masofa $DC=2e$:
\[S_{CFG}=\frac12\cdot\frac h6\cdot2e=\frac{eh}6,\qquad S_{AGE}=\frac12\cdot e\cdot\frac h3=\frac{eh}6.\]
Nisbat $1$.
\Javob{a) 24; \quad b) 1}

\newpage
\Savol{44}
\fig{\begin{tikzpicture}[scale=0.45,font=\small]
\coordinate (A) at (-5.196,0); \coordinate (B) at (5.196,0); \coordinate (C) at (1.732,6); \coordinate (D) at (-1.732,6); \coordinate (O) at (0,3);
\draw (A)--(B)--(C)--(D)--cycle; \fill[okgreen!30] (O) circle (3); \draw (O) circle (3);
\draw[ax] (0,0)--(0,6); \draw[axs] (O)--(3,3); \draw[ax] (C)--(1.732,0);
\fill (O) circle (3pt); \node[left] at (O) {$O$};
\node[below] at (A) {$A$}; \node[below] at (B) {$B$}; \node[above right] at (C) {$C$}; \node[above left] at (D) {$D$};
\node[aux,above] at (1.7,3) {$r$}; \node[aux,below] at (2.6,0) {$R$}; \node[aux,above] at (0.85,6) {$\rho$};
\draw (4.2,0) arc (180:120:1); \node at (3.6,0.55) {\scriptsize$60^\circ$};
\end{tikzpicture}}
Öq kesim — teng yonli trapetsiya $ABCD$, unga aylana içki çizilgan. Katta asos radiusi: $\pi R^2=27\pi$, $R=3\sqrt3$.\\
Yasovçi $l$: \ $R-\rho=l\cos60^\circ=\frac l2$; \ içki çizilgan aylana şarti: $2R+2\rho=2l$, ya'ni $R+\rho=l$.
\[R=\frac{3l}4,\quad \rho=\frac l4\ \Rightarrow\ l=\frac{4R}3=4\sqrt3,\quad \rho=\sqrt3.\]

a) $S_{\text{yon}}=\pi(R+\rho)\,l=\pi\cdot4\sqrt3\cdot4\sqrt3=48\pi$.

b) Balandlik $h=l\sin60^\circ=4\sqrt3\cdot\frac{\sqrt3}2=6$, şar radiusi $r=\frac h2=3$:
\[V=\frac43\pi\cdot27=36\pi.\]
\Javob{a) $48\pi$; \quad b) $36\pi$}

\newpage
\Savol{45}
\fig{\begin{tikzpicture}[xscale=0.4,yscale=0.16,font=\small]
\fill[okgreen!30] plot[domain=0:8,samples=40] (\x,{16-(\x-4)^2})--cycle;
\draw (-1,0)--(11,0); \draw[thick,domain=0:8,samples=40] plot (\x,{16-(\x-4)^2});
\draw (0,0)--(0,36); \draw[axs] (0,36)--(9,0); \draw (4,0)--(4,16); \draw[ax] (0,12)--(6,12);
\node[above] at (0,36) {$A$}; \node[below] at (0,0) {$B$}; \node[below] at (4,0) {$G$}; \node[below left=-1pt] at (8,0) {$E$}; \node[below right=-1pt] at (9,0) {$D$};
\node[above left] at (4,16) {$C$}; \node[left] at (0,12) {$L$}; \node[right] at (6,12.5) {$K$};
\end{tikzpicture}}
$BE=8$, $G$ — $BE$ örtasi, $BG=4$. Parabola: $B$ va $E$ ildizlar, uçi $C$, $CG=16$:
$B$ ni koordinata boşi deb olsak (parabola — funksiya grafigi): \ $y=16-(x-4)^2$.
$K$ nuqta: $KL=6$ $\Rightarrow$ $x=6$, \ $y=16-4=12$, ya'ni $BL=12$.\\
Urinma burçak koeffitsiyenti: $y'=-2(x-4)$, \ $y'(6)=-4$.

a) $AD$ urinma $BE$ bilan $\angle ADE$ hosil qiladi: $\tg\angle ADE=|y'(6)|=4$.

b) $\triangle ALK$ da $\tg\angle AKL=\tg\angle ADE=4$ ($LK\parallel BD$): \ $AL=KL\cdot4=24$.
\[AB=AL+LB=24+12=36.\]
\Javob{a) 4; \quad b) 36}

\ShowJavoblar
\end{document}
